\section{Ruta futura: Generalist Agent}

\subsection{Un modelo omni-modal de tres entradas y tres salidas}

El objetivo final de Qwen es un modelo unificado capaz de \textbf{entender y generar} simultáneamente texto, imagen y audio:

\begin{itemize}
    \item Ya logrado: entendimiento/generación de texto, entendimiento de imagen, entendimiento/generación de audio
    \item Pendiente: incorporar la generación de imagen al modelo unificado
    \item Forma final: «tres entradas y tres salidas»: entrada y salida libres en las tres modalidades
\end{itemize}

\subsection{Multi-turn RL with Environment Feedback}

\begin{importantbox}{Long Horizon Reasoning}
Lin Junyang (林俊旸) sostiene que el futuro del RL no es el autodiálogo al estilo de las competencias matemáticas de DeepSeek R1, sino \textbf{Multi-turn RL with Environment Feedback for Long Horizon Reasoning}:
\begin{itemize}
    \item El modelo ejecuta tareas de cadena larga en un entorno real (por ejemplo, completar en dos días un trabajo que a una persona le tomaría dos meses)
    \item Aprende mediante la retroalimentación del entorno, no solo del autodiálogo
    \item Aplica tanto al mundo virtual (Digital Agent) como al mundo físico (Embodied Agent)
    \item En esencia, convierte una Natural Language Instruction en una Executable Action
\end{itemize}
\end{importantbox}

\begin{knowledgebox}{La tendencia del RL Compute}
La publicidad de XAI muestra que su cómputo de RL ya se acerca al cómputo de Pre-training. Aunque Lin Junyang considera que esto es «un poco un desperdicio», sí demuestra que el RL tiene un enorme margen de posibilidades. Al equipo de Qwen no le interesa tanto que el modelo sea el cerebro matemático más poderoso, sino que pueda contribuir a la biología y a la sociedad como lo haría una persona real.
\end{knowledgebox}

\subsection{Del Digital Agent al Physical Agent}

\begin{itemize}
    \item \textbf{Digital Agent}: opera simultáneamente la GUI y la API; es la dirección más madura por ahora
    \item \textbf{Physical Agent}: realiza tareas del mundo físico como manejar un micrófono o servir agua; requiere el soporte de VLA (Vision-Language-Action)
    \item Perspectiva unificada: tanto el Coding Agent como el VLA consisten, en esencia, en «convertir instrucciones en lenguaje natural en acciones ejecutables»
\end{itemize}

\subsection{Resumen de la sección}

La ruta técnica de Qwen es clara: unificación omni-modal → Agent RL → Digital/Physical Agent. Es un camino de evolución que va de «entender el mundo» a «cambiar el mundo».

